\section{The exact regrouping identity}
\label{sec:ope-grouping-algebra}
The estimates of Section~\ref{sec:ope-grouping} concern an exact
reorganization of the marked density. We now prove the identity. In
particular the coefficient of the empty mark set, which determines
the vacuum activity, is preserved.

\subsection{Squarefree marks and complete groups}
For a finite mark set $I$, use the commutative algebra with basis
$e_A$, $A\subseteq I$, and multiplication
\[
 e_Ae_B=\begin{cases}e_{A\cup B},&A\cap B=\varnothing,\\0,&\text{otherwise}.
 \end{cases}
\]
Its unit is $e_\varnothing$. The ideal spanned by the nonempty
$A$ is nilpotent, so its exponential is a finite sum. In particular,
for finitely many nonempty $A_r$,
\[
 \exp\left(\sum_r z_re_{A_r}\right)
       =\prod_r(1+z_re_{A_r}),
\]
since each individual summand squares to zero and all commute.
Repeated use of a mark thus vanishes automatically.

Let $\pi:I\to V$ be the partition from the previous section, with
nonempty groups $I_v=\pi^{-1}(v)$. Define a linear map
$\mathcal R$ to the analogous algebra on $V$ by
\[
 \mathcal R e_A=
 \begin{cases}e_B,&A=\bigcup_{v\in B}I_v\text{ for some }B\subseteq V,\\
 0,&\text{otherwise}.
 \end{cases}
\]
Equivalently, the coefficient at $B$ after applying $\mathcal R$
is the old coefficient at $\pi^{-1}(B)$. This map is not asserted
to preserve arbitrary products. It does preserve a product of
monomials whose mark sets are individually unions of complete groups.
If these sets are disjoint, the assertion follows from their unions;
if they overlap, both products vanish. We will apply multiplicativity
only after the connected-component decomposition has made every
component complete in this sense.

\subsection{Ordinary and flagged entries}
Fix a lattice configuration and let $F$ be its set of flags. Take
finite families of ordinary and flagged entries. An
ordinary entry $r$ has a nonempty old mark set $A_r$, a support
$X_r$ and a coefficient $z_r$. A flagged entry $r$ has, in addition,
a nonempty prescribed flag set $Z_r\subseteq X_r$, and its coefficient
vanishes unless $X_r\cap F=Z_r$. Its mark set may be empty.
The flagged part of the density is the sum over selections of flagged
entries with pairwise disjoint supports and prescribed flag sets whose
union is $F$. Multiply this sum by the exponential of the ordinary
entries. The finite exponential formula shows that the product is a
sum over sets $S$ of distinct ordinary and flagged entries, with weight
\[
 \prod_{r\in S}z_r e_{A_r},
\]
subject to the flagged-support and flag-inventory conditions. Terms
with overlapping old mark sets vanish.

Enlarge each support as in \eqref{eq:ope-enlarged-cover}, writing
$\widetilde X_r$ for the result. Split each ordinary entry into two
entries with the same marks, multiplied respectively by
$\boldsymbol1_{\widetilde X_r\cap F=\varnothing}$ and its complement.
Call the second kind \emph{dirty}; all flagged entries are dirty too.
The sum of the two pieces is the original entry, so this split leaves
the density unchanged. At most one piece is nonzero at any fixed
configuration. Every enlarged support contains the representative
$c_v$ of each group mentioned by its old marks. A flagged entry with
no marks is not enlarged.

We describe a canonical partition of a selected set $S$. First join
entries which mention a common group. Call a component of this first
graph hot if it contains a dirty entry. Next join hot components
whenever their enlarged supports intersect. The resulting components
are the components used below. Equivalently each entry always carries
its group names as graph coordinates, and carries its enlarged support
as additional graph coordinates precisely when its first component
is hot. Connected components are then taken in this single intersection
graph. Distinct components have disjoint group-name sets, and distinct
hot components have disjoint enlarged supports. Every component is
connected by ordinary support overlaps: a common group supplies its
representative site as a common point.

This partition is unchanged by restricting to one of its components.
Indeed a path using group names never leaves that component, so its
hot status can be determined within it. Conversely, a family of
internally connected components with disjoint group-name sets and
pairwise disjoint hot supports is exactly the component partition
of its union. No group-name path or hot-support edge can cross between
members of the family. These two observations prove the uniqueness
needed for the reorganization.

\subsection{Reconstructing the selected density}
Consider a nonzero old term which survives $\mathcal R$. Its total
mark set is a union of complete groups. Every one of its components
also has this property. For if one mark in a group occurs in that
component, every other selected mark of the same group is joined to
it through their common group name. Since the total selection contains
the whole group, the component contains the whole group.

A cold component contains only healthy ordinary entries, so its
coefficient is supported on configurations with no flags in its
combined enlarged support. It has a nonempty group mark set. A hot
component necessarily contains a flagged entry. To prove this last
assertion, choose a dirty ordinary entry if no flagged entry has yet
been found. Its enlarged support meets $F$. The global flag-inventory
condition supplies a selected flagged entry whose prescribed flag set
contains that point. Both entries are hot and their supports intersect,
so they are in the same component, a contradiction.

For each possible cold component, make a new ordinary entry: its
support is the union of its enlarged supports, its coefficient is
the product of its constituent coefficients, and its new marks are
its group names. For each possible component containing a flagged
entry, make a new flagged entry with that same combined support,
flag set $Z$ equal to the union of its old prescribed flag sets,
and coefficient
\[
 \boldsymbol1_{\widetilde X\cap F=Z}\prod_{r\in S}z_r.
\]
Only internally connected components with disjoint old mark sets,
complete group marks and pairwise disjoint old flagged supports are
included. The new flag set is nonempty.

For the components of a nonzero old selection, the displayed indicator
is one. Any flag in the enlarged support of a hot component belongs
to a selected flagged entry somewhere in the old term; intersection
of the two hot supports puts that entry in this same component.
The converse inclusion follows from the old flag conditions.
Distinct hot supports are disjoint, and their prescribed flag sets
still partition $F$. Distinct components have disjoint new mark sets.
The old term therefore gives an allowed new term with exactly the
same coefficient after $\mathcal R$.

Conversely, take a nonzero term in the new ordinary-exponential and
flagged-selection expansion. Its component entries have disjoint
new mark sets and pairwise disjoint hot supports. Their inventories
of old labels are disjoint. A shared old label carrying a mark would
force a shared group name; a shared label without marks must be a
flagged entry and would force intersecting hot supports. Thus no old
label is counted twice. Within each new component the old flagged
supports are disjoint; between components they lie inside disjoint
hot supports. Their old flag sets have union $F$. The old mark sets
are disjoint, since they are disjoint within components and complete
disjoint groups between components. Hence their union gives an
allowed old selection. The uniqueness of the component partition
proves that this construction inverts the preceding one.

The products of coefficients agree in both directions. On complete
group components, $\mathcal R$ preserves their monomial product by
the first subsection. Summing the resulting weight-preserving
bijection proves the exact regrouping identity. All sums here are
finite: the torus, mark sets, supports and their admissible patterns
are finite. No convergence or rearrangement assumption enters this
algebraic step. The connected-array estimates of
Section~\ref{sec:ope-grouping} apply to these very component entries.
The ordinary and flagged kinds cannot represent the same constituent
set, since one is cold and the other contains a flagged entry; thus
summing their bounds does not hide an additional multiplicity.

\subsection{Vacuum, locality, and covariance}
A new entry with no group marks must be a single old zero-mark flagged
entry. In fact no ordinary entry has empty marks. All entries of such
a component are therefore flagged entries with no marks, whose
supports were not enlarged. Their old supports are pairwise disjoint,
so no support edge connects two different entries; and there are no
group-name edges. Connectedness forces a singleton. Its new support,
flag set and coefficient coincide with the old ones; the indicator
is already implied by the old flag condition. Thus the zero-mark
coefficient of every flagged support is unchanged. The ordinary
array has no zero-mark coefficient. The same vacuum activity is
retained by the regrouping.

The construction also preserves the required local dependence.
Each coefficient is a finite product of original coefficients, with
indicators determined by the flags in the combined support. It
therefore depends only on the original field and flag data in that
support. Products, finite sums and these measurable indicators preserve
measurability. If the original coefficients and flags are gauge
invariant, so are the new coefficients. Healthy ordinary components
vanish when their support meets the flags; flagged components vanish
unless the intersection with the flags is their specified flag set.
These are the support conditions used in the local representations.
They do not assert smoothness across a flag boundary.

Finally we state covariance. Let $a$ be a permutation of the old marks
and $b$ a permutation of the group names, and let a translation or signed
coordinate change, with the accompanying conjugation, act on positions
and flags. Assume: (a) $a$ and $b$ are compatible with the partition,
$\pi(a(i))=b(\pi(i))$ for every old mark $i$, and they carry the chosen
representative of each group to the chosen representative of its image;
(b) for a fixed copied family $v$ of group names with image $u$, one has
$b(i)\in u$ if and only if $i\in v$; (c) for marks in groups of $v$, the
positions and flag data of $a(i)$ are the transformed positions and flag
data of $i$; and (d) the original coefficient identities hold for every
mark set contained in the copied family. Compatibility (a) is an
assumption on the relabelling, not a consequence of (d): a permutation of
old marks that does not map groups to groups does not intertwine the
regrouping. In the application the same partition is used in each copy,
as in Section~\ref{sec:ope-cutoff}, and (a)--(c) hold by construction.
Under (a)--(d) the transformation relabels entries bijectively, preserves
group-name incidence, health and flag inventory, and sends support
intersections to support intersections. Hence they send the canonical
components to the corresponding components after transformation.
Reindexing the finite sums proves covariance of the new arrays for
every coefficient whose group names lie in the copied family $v$; no
statement is made for other coefficients. This keeps the replication identities through repeated
regrouping. To obtain the local normalized-field representation used
in \eqref{eq:ope-group-weight}, one must in addition extract its
singleton coefficients and apply the purchased local analytic
projection bounds; those estimates are not consequences of this
finite algebraic identity.
